\chapter{CALL: Automated Feedback Classification Coding Schemes}\label{app:call_coding}
\begin{CJK}{UTF8}{min}

All feedback messages generated by the chatbot were classified across three feedback models (Explicitness, Directness, and Writing Skill) using the LLM classification pipeline.

\section{Classification Prompt Template}

Each category was evaluated independently using the following prompt structure:

\medskip\hrule\medskip

\# MASTER PROMPT

\medskip\hrule\medskip

{[}System{]}

You are an expert EFL writing feedback researcher. Decide if the chatbot's message contains feedback targeting \{CATEGORY\_NAME\}.

Definition: \{CATEGORY\_DEFINITION\}

Inclusion Criteria (True): \{INCLUSION\_CRITERIA\}

Exclusion Criteria / Guardrails (False): \{EXCLUSION\_CRITERIA\}

Respond with JSON: \{"rationale": string, "present": boolean\}

{[}User{]}

Chatbot message:

"""

\{penny\_message\}

"""

Is this feedback category present? Reason step by step, then give your answer in JSON.

\section{Explicitness}

\medskip\hrule\medskip

\# EXPLICIT CORRECTION PROMPT

\medskip\hrule\medskip

"""You are an expert EFL writing feedback researcher. Decide if a chatbot's message contains **Explicit Correction** feedback.

DEFINITION

Explicit Correction is when the teacher/chatbot clearly and explicitly names/states that a specific word, phrase, or sentence was changed and provides the specific corrected form. It explicitly tells the student what was fixed.

TRUE EXAMPLES:

- "「then」は「them」に修正しました。" (I corrected 'then'to 'them'.)

- "「and」を追加して、最後の文をつなげてみました。" (I added 'and'to connect the last sentences.)

- "Change 'houseworks'to 'housework'."

FALSE EXAMPLES (e.g., pure recast without pointing out the change):

- "Here is the revised text: {[}Corrected Text{]}." (If it doesn't explicitly name the specific change made).

"""

\medskip\hrule\medskip

\# RECAST PROMPT

\medskip\hrule\medskip

"""You are an expert EFL writing feedback researcher. Decide if a chatbot's message contains **Recast** feedback.

DEFINITION

Recast is when the teacher/chatbot provides a reformulated, corrected version of the student's entire utterance or a substantial part of it, without necessarily stopping the flow to point out the specific error. In chatbot logs, this usually appears as a "Revised text:" block.

TRUE EXAMPLES:

- "**Revised text:** If I could meet any historical figure, I would want to meet..."

- "**修正後の文章:** I want to have dinner with Sugino Yosuke..."

FALSE EXAMPLES:

- Just asking a clarification question.

- Pointing out a typo without providing a rewritten chunk of the user's sentence.

"""

\medskip\hrule\medskip

\# ERROR HIGHLIGHT PROMPT

\medskip\hrule\medskip

"""You are an expert EFL writing feedback researcher. Decide if a chatbot's message contains **Error Highlight** feedback.

DEFINITION

Error Highlight occurs when the teacher/chatbot emphasises the exact location of an error or the specific correction within a sentence, but DOES NOT provide an explicit correction statement.

CRITICAL GUARDRAILS (MARK FALSE IF THESE OCCUR):

\textbf{[NO]} FALSE: If the ONLY bold text is a header/title like "**Revised text:**" or "**修正後の文章:**".

\textbf{[NO]} FALSE: If the chatbot bolds an ENTIRE paragraph or ENTIRE sentence (e.g., "**I went to the store and bought an apple.**"). This is not highlighting a specific error; it is just formatting a block of text.

TRUE EXAMPLES (MARK TRUE ONLY IF):

\textbf{[YES]} TRUE: Bolding is used inline on a specific word/phrase to show contrast inside a normal sentence or bulleted list to draw attention to the exact location of the fix. Explicit correction statements are not present, but the bolding serves to highlight the error or the correction.

Example: "1. went **to** Dorogawa hot spring". How can you improve this sentence? The error is highlighted in bold.

Example: "I visited a convenience store **called RO,SON**." Correct the error in this sentence.

"""

\medskip\hrule\medskip

\# CLARIFICATION REQUEST PROMPT

\medskip\hrule\medskip

"""You are an expert EFL writing feedback researcher. Decide if a chatbot's message contains **Clarification Request** feedback.

DEFINITION

A Clarification Request is when the chatbot explicitly asks the student to explain their intended meaning because the chatbot genuinely does not understand what the student is trying to say in their original English text.

CRITICAL GUARDRAILS (MARK FALSE IF THESE OCCUR):

\textbf{[NO]} FALSE: Administrative questions about what to do next: "次は、どの部分にもっと集中したいですか？" (What part do you want to focus on next?) -\textgreater{} FALSE

\textbf{[NO]} FALSE: Prompts to write in the box: "テキストボックスに書き直してみてください。" (Please rewrite it in the text box) -\textgreater{} FALSE

\textbf{[NO]} FALSE: Conversational closers: "他に質問がありますか？" (Do you have any other questions?) -\textgreater{} FALSE

\textbf{[NO]} FALSE: Asking for confirmation: "これで大丈夫ですか？" (Is this okay?) -\textgreater{} FALSE

\textbf{[NO]} FALSE: Conversational elicitation or follow-up questions meant to encourage more writing, rather than clarifying a confusing sentence (e.g., "具体的にどの作品が好きですか？" - What specific work do you like?).

\textbf{[NO]} FALSE: Offering conditional help for future writing (e.g., "もし文を考えている場合は教えてください" - If you are thinking of a sentence, please let me know).

TRUE EXAMPLES (MARK TRUE ONLY IF):

\textbf{[YES]} TRUE: "具体的にどのような使い方を考えているのか教えていただけますか？" (Could you tell me specifically how you want to use it / what you mean?)

\textbf{[YES]} TRUE: "「どろどろ」は文脈によって意味が異なります。どのような意味で使いたいですか？" (Depending on the context, 'dorodoro'means different things. What meaning do you want to use?)

"""

\medskip\hrule\medskip

\# METALINGUISTIC FEEDBACK PROMPT

\medskip\hrule\medskip

"""You are an expert EFL writing feedback researcher. Decide if a chatbot's message contains **Metalinguistic** feedback.

DEFINITION

Metalinguistic feedback is present only when the chatbot explicitly comments on the learner's language form or use and explains the linguistic reason for a change, such as grammar, syntax, tense, article use, plurality, spelling, word choice, register, discourse markers, or naturalness. It may name the error type or explain the rule/nuance behind the correction.

CRITICAL GUARDRAILS (MARK FALSE IF THESE OCCUR):

\textbf{[NO]} FALSE: Explaining the definition of a Japanese word the student asked about (e.g., defining what "どろどろ" means), rather than explaining a linguistic rule about their English text.

\textbf{[NO]} FALSE: Stating a mechanical change without linguistic reasoning, such as "「then」は「them」に修正しました。" or "I changed X to Y."

\textbf{[NO]} FALSE: Giving only a reformulated answer, praise, or a vague comment like "this sounds better", "more natural", "more fluent", or "better flow" without naming the language issue.

\textbf{[NO]} FALSE: Asking for clarification or inviting the student to continue without explaining the error.

\textbf{[NO]} FALSE: Explaining the revision only as a stylistic preference or general improvement unless the message also names the language issue itself.

TRUE EXAMPLES (MARK TRUE ONLY IF):

\textbf{[YES]} TRUE: Providing a grammatical label or categorising the error type and linking it to the learner's text (e.g., "(単数形に修正)", "(冠詞を追加)", "(スペルミス)").

\textbf{[YES]} TRUE: Explains a semantic reason for a change that refers to the learner's wording (e.g., "「nothing something」は不明瞭なので削除しました").

\textbf{[YES]} TRUE: Explains nuance such as formality, politeness, or naturalness in a way that makes the linguistic reason explicit: "「have dinner」を「dine」に変えて、少しフォーマルにしてみました。"

\textbf{[YES]} TRUE: Mentions a rule or form issue directly, such as "You need the past tense here" or "Use the singular form here."

\textbf{[YES]} TRUE: Explains a discourse-marker or connective choice, such as why "Additionally" is used to add information or why "another" is different from "other".

\textbf{[YES]} TRUE: Explains a contrast in meaning or usage, such as "near" vs. "close to" or a singular/plural reference choice.

DECISION RULE

If the message only shows the correction or summarises the edit, mark FALSE.

If the message says the wording is "better", "more natural", or "smoother" but does not explain the actual language issue, mark FALSE.

If the message names the language issue, rule, or nuance that motivated the correction, mark TRUE.

"""

\section{Directness}

\medskip\hrule\medskip

\# DIRECTIVE PROMPT

\medskip\hrule\medskip

"""You are an expert EFL writing feedback researcher. Decide if a chatbot message contains **Directive** feedback.

DEFINITION

Directive feedback tells the learner exactly what to change, add, remove, or rewrite.

MARK TRUE when the message includes any of these:

- A corrected rewrite or revised version of the learner's text.

- A specific replacement, deletion, insertion, or edit, such as "X -\textgreater{} Y" or "change X to Y".

- A direct instruction to modify specific text, grammar, wording, organisation, or content.

MARK FALSE when the message only gives praise, asks what to do next, explains a concept, or gives vague advice without a specific fix.

Respond with JSON: \{response\_format\}

"""

\medskip\hrule\medskip

\# INFORMATIVE PROMPT

\medskip\hrule\medskip

"""You are an expert EFL writing feedback researcher. Decide if a chatbot message contains **Informative** feedback.

DEFINITION

Informative feedback gives the learner useful knowledge, content guidance, or conceptual commentary for revision. It tells the learner what to think about or understand, not merely which exact correction was made.

MARK TRUE when the message:

- Gives content-development advice, such as adding topics, details, reasons, examples, hobbies, goals, or explanations.

- Names a broad improvement dimension without an exact correction, such as making sentences clearer, using more specific words, organising information, or adding detail.

- Explains a concept or meaning the learner can use in the writing, such as what "黒歴史" means or how it relates to the topic.

- Simplifies a difficult idea, explains a term, or offers multiple possible expressions so the learner can choose what fits.

MARK FALSE when the message is only:

- A corrected rewrite or revised text, even if it says "pay attention to grammar/vocabulary".

- A full example passage, model sentence, translation, or "write it like this" text that the learner is expected to copy or reflect directly into the revision.

- A specific edit or correction note, such as "X -\textgreater{} Y", "I changed X to Y", "add the", "use plural", or "use past tense".

- A short reason attached to a specific correction, such as "because the subject is plural", "this is more natural", "this is more formal", "this makes the sentence smoother", or "the wording is unclear".

- A list of correction categories after a rewrite, such as "grammar, spelling, vocabulary, tense".

- One exact translation or expression followed by "use this expression", unless the message also explains the term/concept or gives multiple alternatives.

- A genuine clarification question asking what the learner means or what topic they are writing about; that is Query, not Informative.

- Praise, greeting, or administrative follow-up.

Decision rule: If the message mainly gives text to copy, explains corrections already made in a revised text, or asks for clarification, mark FALSE. If it gives broader content guidance, conceptual explanation, simplification, or choices the learner can apply beyond one exact correction, mark TRUE.

"""

\medskip\hrule\medskip

\# QUERY PROMPT

\medskip\hrule\medskip

"""You are an expert EFL writing feedback researcher. Decide if a chatbot message contains **Query** feedback.

DEFINITION

Query feedback is a genuine question asking the learner to clarify their intended meaning, content, or direction in the writing.

MARK TRUE only when the chatbot is genuinely uncertain about what the learner means and asks for clarification.

MARK FALSE for next-step questions, revision prompts, translation offers, reflection questions, or general conversation such as "What part would you like to improve next?" or "Please show me the revised text."

"""

\medskip\hrule\medskip

\# PRAISE PROMPT

\medskip\hrule\medskip

"""You are an expert EFL writing feedback researcher. Decide if a chatbot message contains **Praise** feedback.

DEFINITION

Praise feedback positively evaluates or encourages the learner's writing, effort, or improvement.

MARK TRUE for direct compliments such as "great job", "your writing is very good", "well done", or positive comments about a specific improvement.

MARK FALSE for polite openers like "Of course" or "Sure", and for messages that are only corrections or administrative follow-up.

"""

\section{Feedback Focus}

\medskip\hrule\medskip

\# SPELLING PROMPT

\medskip\hrule\medskip

"""You are an expert writing feedback researcher. Decide if a chatbot's message contains feedback specifically targeting **Spelling**.

DEFINITION:

Feedback that addresses the correct arrangement of letters in a word.

"""

\medskip\hrule\medskip

\# PUNCTUATION \& CAPITALISATION PROMPT

\medskip\hrule\medskip

"""You are an expert writing feedback researcher. Decide if a chatbot's message contains feedback specifically targeting **Punctuation** or **Capitalisation**.

DEFINITION:

Feedback that addresses the use of standard marks (commas, periods, quotation marks, exclamation points, etc.) to separate words or sentences.

"""

\medskip\hrule\medskip

\# GRAMMAR PROMPT

\medskip\hrule\medskip

"""You are an expert writing feedback researcher. Decide if a chatbot's message contains feedback specifically targeting **Grammar**.

DEFINITION:

Feedback that addresses grammatical rules, such as subject-verb agreement, verb tense, article usage, prepositions, plurals, or pronoun references etc. Only consider grammar feedback that is explicitly mentioned in the chatbot's message. I.e. not just a recasting of the student's writing, but a direct comment on the grammar.

"""

\medskip\hrule\medskip

\# WORD CHOICE PROMPT

\medskip\hrule\medskip

"""You are an expert writing feedback researcher. Decide if a chatbot's message contains feedback specifically targeting **Word Choice**.

DEFINITION:

Feedback that addresses the specific vocabulary selected by the student. This includes suggesting more precise vocabulary, sophisticated phrasing, or pointing out misused words. Only consider words that are explicitly mentioned in the feedback message. I.e. not just a recasting of the student's writing, but a direct comment on the words themselves.

"""

\medskip\hrule\medskip

\# ORGANISATION PROMPT

\medskip\hrule\medskip

"""You are an expert writing feedback researcher. Decide if a chatbot's message contains feedback specifically targeting **Organisation**.

DEFINITION:

Feedback that addresses the structure and flow of the writing. This includes suggestions for improving transitions, paragraphing, and overall coherence.

"""

\medskip\hrule\medskip

\# IDEAS \& ELABORATION PROMPT

\medskip\hrule\medskip

"""You are an expert writing feedback researcher. Decide if a chatbot's message contains feedback specifically targeting **Ideas and Elaboration** (a higher-level writing skill).

DEFINITION:

Feedback directed at the content of the writing. This includes prompting the user to add more sensory details, expanding on a thought, answering queries about the narrative (e.g., "What did it look like?"), or improving the overall concepts and structures in the text.

"""
\end{CJK}

